Next, to verify whether these characteristics are unique to highly cited papers, Table \ref{tab:g2obsindex} presents the index values for G2 papers.
Since calculating these index values requires a sufficient number of citations, we limit the comparison group to G2 papers only.
According to these tables, highly cited papers (G1) are generally older than G2.
In terms of class composition, highly cited papers show a significantly higher proportion of Class G and Class SB compared to G2, indicating a clear difference.
Furthermore, we examined the peak year (the median time from publication to peak citation) for the group that has sufficiently matured and exhibits an obsolescence tendency, i.e., Class No. Comparing G1 and G2, the values were 9.5 years and 5 years, respectively. Compared to existing reports (e.g., Parolo et al.), G1 tended to show longer durations, whereas G2 was consistent with previous findings.

